%------------------------------------------------------------------------------%
\begin{question}
  Resolva o sistema, utilizando escalonamento,
  % \[
  %   \systeme{
  %      x +  y + z = 6,
  %     2x -  y + z = 3,
  %      x + 2y - z = 2}
  % \]
  \[
    \left\{\begin{array}{rrrr}
      x  & +  y & + z & = 6 \\
      2x & -  y & + z & = 3 \\
      x  & + 2y & - z & = 2
    \end{array}\right.
  \]
\end{question}

%------------------------------------------------------------------------------%
\begin{resolution}{Matriz Aumentada}
  Matriz aumentada
  \[
    \begin{amatrix}{3}
      1 & 1  & 1  & 6 \\
      2 & -1 & 1  & 3 \\
      1 & 2  & -1 & 2
    \end{amatrix}
  \]
\end{resolution}

%------------------------------------------------------------------------------%
\begin{resolution}{Escalonamento do Sistema}
\begin{columns}
\column{0.4\linewidth}
  \[\;\;\,
    \begin{amatrix}{3}
      1 & 1  & 1  & 6 \\
      2 & -1 & 1  & 3 \\
      1 & 2  & -1 & 2
    \end{amatrix}
  \]
  \pause
\column{0.6\linewidth}
  Eliminar os elementos abaixo do primeiro pivô
  \\ \pause
  \(
    L_2 \leftarrow L_2 - 2 L_1
  \)
  \\ \pause
  \(
    L_3 \leftarrow L_3 - L_1
  \)
\end{columns}
\vspace{-\baselineskip}
\[
  \pause
  L'_2
  \;=\;
  L_2 - 2 L_1
  \pause
  \;=\;
  \begin{amatrix}{3}
    2 & -1 & 1  & 3
  \end{amatrix}
  - 2
  \begin{amatrix}{3}
    1 & 1  & 1  & 6
  \end{amatrix}
  \pause
  \;=\;
  \begin{amatrix}{3}
    0 & -3 & -1 & -9
  \end{amatrix}
\]
\[
  \pause
  L'_3
  \;=\;
    L_3 - L_1
  \pause
  \;=\;
  \begin{amatrix}{3}
    1 & 2  & -1 & 2
  \end{amatrix}
  -
  \begin{amatrix}{3}
    1 & 1  & 1  & 6
  \end{amatrix}
  \pause
  \;=\;
  \begin{amatrix}{3}
    0 & 1  & -2 & -4
  \end{amatrix}
\]
\pause
\[
  \begin{amatrix}{3}
    1 & 1  & 1  & 6  \\
    0 & -3 & -1 & -9 \\
    0 & 1  & -2 & -4
  \end{amatrix}
\]
\end{resolution}

%------------------------------------------------------------------------------%
\begin{resolution}{Escalonamento do Sistema}
\begin{columns}
\column{0.4\linewidth}
  \[\;\;\,
  \begin{amatrix}{3}
      1 & 1  & 1  & 6  \\
      0 & -3 & -1 & -9 \\
      0 & 1  & -2 & -4
  \end{amatrix}
  \]
  \pause
\column{0.6\linewidth}
  Trocamos a segunda e terceira linhas
  \\
  \(
    L_2\leftrightarrow L_3
  \)
\end{columns}
\vspace{-\baselineskip}
\[
  \pause
  L'_2
  \;=\;
  L_3
  \pause
  \;=\;
  \begin{amatrix}{3}
    0 & 1  & -2 & -4
  \end{amatrix}
\]
\[
  \pause
  L'_3
  \;=\;
  L_2
  \pause
  \;=\;
  \begin{amatrix}{3}
    0 & -3 & -1 & -9
  \end{amatrix}
\]
\pause
\[
  \begin{amatrix}{3}
    1 & 1  & 1  & 6  \\
    0 & 1  & -2 & -4 \\
    0 & -3 & -1 & -9
  \end{amatrix}
\]
\end{resolution}

%------------------------------------------------------------------------------%
\begin{resolution}{Escalonamento do Sistema}
\begin{columns}
\column{0.4\linewidth}
  \[\;\;\,
  \begin{amatrix}{3}
      1 & 1  & 1  & 6  \\
      0 & 1  & -2 & -4 \\
      0 & -3 & -1 & -9
  \end{amatrix}
  \]
  \pause
\column{0.6\linewidth}
  Eliminar os elementos abaixo do segundo pivô
  \\ \pause
  \(
    L_3 \leftarrow L_3 + 3 L_2
  \)
\end{columns}
\vspace{-\baselineskip}
\[
  \pause
  L'_3
  \;=\;
    L_3 + 3 L_2
  \pause
  \;=\;
  \begin{amatrix}{3}
    0 & -3 & -1 & -9
  \end{amatrix}
  + 3
  \begin{amatrix}{3}
    0 & 1  & -2 & -4
  \end{amatrix}
\]
\[
  \pause
  \hphantom{L'_3}
  \;=\;
  \begin{amatrix}{3}
    0 & 0 & -7 & -21
  \end{amatrix}
\]
\pause
\[
  \begin{amatrix}{3}
    1 & 1 & 1  & 6   \\
    0 & 1 & -2 & -4  \\
    0 & 0 & -7 & -21
  \end{amatrix}
\]
\end{resolution}

%------------------------------------------------------------------------------%
\begin{resolution}{Sistema Linear Escalonado}
  Matriz aumentada do sistema escalonado
  \[
    \begin{amatrix}{3}
      1  & 1  & 1  & 6   \\
      \z & 1  & -2 & -4  \\
      \z & \z & -7 & -21
    \end{amatrix}
  \]
  
  \pause
  Sistema linear escalonado
  % \[
  %   \systeme{
  %     x + y +  z =  6,
  %         y - 2z = -4,
  %           - 7z = -21}
  % \]
  \[
    \left\{\begin{array}{rrrr}
      x & + y & +  z & =  6 \\
        &   y & - 2z & = -4 \\
        &     & - 7z & = -21
    \end{array}\right.
  \]
\end{resolution}

%------------------------------------------------------------------------------%
\begin{resolution}{Substituição Retroativa}
 \begin{columns}[t]
 \column{0.26\linewidth}
 \onslide<+->{\centerline{Última linha}}
 \begin{align*}
   \onslide<+->{-7 z & = -21} \\
   \onslide<+->{z    & = 3}
 \end{align*}
 \column{0.27\linewidth}
 \onslide<+->{\centerline{Segunda linha}}
 \begin{align*}
   \onslide<+->{y -2 z        & = -4} \\
   \onslide<+->{y -2 \times 3 & = -4} \\
   \onslide<+->{y -6          & = -4} \\
   \onslide<+->{y             & = 2}
 \end{align*}
 \column{0.4\linewidth}
 \onslide<+->{\centerline{Primeira linha}}
 \begin{align*}
   \onslide<+->{x + y + z & = 6} \\
   \onslide<+->{x + 2 + 3 & = 6} \\
   \onslide<+->{x         & = 1}
 \end{align*}
 \end{columns}
\end{resolution}

%------------------------------------------------------------------------------%
\begin{resolution}{Solução}
A solução do sistema é
\(
  \mathbf{x}
  \;=\;
  \begin{bmatrix}
    x \\ y \\ z
  \end{bmatrix}
  \;=\;
  \begin{bmatrix}
    1 \\ 2 \\ 3
  \end{bmatrix}
\)
\end{resolution}

%------------------------------------------------------------------------------%
